
\subsubsection{6.1 Key Interpretations}

\textbf{The Checkbox Compliance Pattern:} Our most important finding is not the null result (p=0.418) but the reason for it: artifact quality is critically low (2.43/10). This pattern is consistent with checkbox compliance: reproducibility badge programs have succeeded at increasing artifact presence but not artifact quality. Authors create GitHub repositories and deposit minimal documentation to satisfy venue requirements, but they do not invest time in detailed specifications. The result is a proliferation of badge-compliant artifacts that provide little replication value. Alternative explanations include time constraints, venue enforcement gaps, and inadequate tooling support.

\textbf{Why Evaluation Protocols Are Missing:} The near-zero scores for evaluation protocols (1.19/10) and hyperparameters (1.16/10) are particularly striking. These dimensions are straightforward to document---evaluation scripts can be shared verbatim, hyperparameters fit in a single table---yet they are systematically absent. We hypothesize two contributing factors: (1) implicit knowledge assumption: authors assume evaluation protocols are obvious for standard benchmarks; (2) post-publication neglect: authors create artifacts during the initial submission but do not maintain them post-acceptance.

\textbf{The Underpowered Trend:} While the Mann-Whitney test was non-significant (p=0.418), the effect size (d=0.464) approaches the medium threshold (0.5), and the directional trend (mean CV: 0.035 vs 0.069) is consistent with our hypothesis. Our sample size (n=22) was far below the target (n=100), resulting in only ~30\% statistical power versus the intended 80\%. This raises the possibility of a Type II error: a real but weak effect exists, but our study was underpowered to detect it.

\subsubsection{6.2 Unexpected Findings and Alternative Explanations}

\textbf{No Dose-Response Relationship:} We expected artifact count (0-3) to correlate negatively with performance variance, but observed $\rho$=-0.084 (negligible). Two competing explanations emerge:
\begin{enumerate}
\item \textbf{Quality dominates quantity:} One high-quality artifact outweighs three low-quality ones. Since h-m1 found most artifacts are low-quality, adding more provides no marginal benefit.
\item \textbf{Confounding by popularity:} Popular benchmarks (ImageNet, CIFAR-10) attract both more artifacts and community-standardized protocols independent of artifact content. Our observational design cannot disentangle these effects.
\end{enumerate}

\textbf{Why ObjectNet Is an Outlier:} ObjectNet showed extreme variance (CV=0.293), inflating the low-artifact group mean. However, this is not a flaw in our data---ObjectNet is designed to test robustness under distribution shift, so high variance reflects its purpose. This illustrates a broader challenge: in observational studies, task characteristics (difficulty, domain, maturity) confound artifact effects.

\subsubsection{6.3 Limitations}

\textbf{Sample Size:} Our variance analysis (h-m3) was severely underpowered (n=22 vs target n=100). Manual data collection from 58 papers was time-intensive, and the Papers with Code API was unavailable during our study period. This limitation has two implications: (1) the non-significant result (p=0.418) could be Type II error---a real d~0.5 effect might exist but be undetectable at n=22; (2) confidence intervals are wide, limiting precision of effect size estimates. We report the effect size (d=0.464) and power analysis, enabling future meta-analyses to incorporate our findings. A directional trend with small sample is more informative than no data.

\textbf{Artifact Quality Measurement:} We used automated content analysis (keyword-based rubric scoring) rather than expert human raters, which may underestimate quality if artifacts use non-standard terminology. However, perfect automated scoring consistency ($\kappa$=1.0) and validation against real artifact content suggest the approach is valid. The binary nature of quality distinctions (complete specification vs minimal information) reduces subjectivity.

\textbf{CV Measures Consistency, Not Correctness:} Our primary metric (coefficient of variation) quantifies procedural consistency across independent attempts, not whether results are correct. Multiple labs could consistently reproduce wrong results if they all inherit the same implementation bug. We measure one necessary condition for reproducibility (consistency) but not sufficiency. Inconsistency is a reproducibility failure regardless of correctness. If labs cannot consistently replicate results, the method is not reproducible.

\textbf{h-m2 Incomplete:} Our planned analysis of protocol consistency (Step 2 of the causal mechanism) was blocked by Semantic Scholar API rate limiting. We cannot directly verify whether low artifact quality (h-m1) translates to high protocol ambiguity. However, the convergence of h-m1 (low quality) and h-m3 (no variance reduction) provides indirect evidence that the mechanism failed at Step 2-3.

\subsubsection{6.4 Connection to Existing Literature}

\textbf{Confirming Prior Barriers Research:} Semmelrock et al. (2024) identified documentation as a reproducibility barrier in their qualitative framework. Our work quantifies the severity: mean quality 2.43/10, with evaluation protocols and hyperparameters almost entirely missing. This transforms a qualitative observation into a quantifiable policy target.

\textbf{Extending Leakage Work:} Kapoor \& Narayanan (2023) showed that data leakage affects hundreds of papers despite documentation requirements. Our finding---that artifacts exist but lack detail---explains their observation: checkbox compliance produces artifacts that satisfy formal requirements but do not prevent methodological errors. Documentation alone is insufficient without quality enforcement.

\textbf{Replicating FAIR Compliance Findings:} Gim et al. (2025) found 5\% FAIR Findable and 0\% Reusable in medical imaging datasets. We replicate this pattern in the ML benchmark domain: artifact presence (Findable) is high, but artifact quality (Reusable) is low. This suggests the problem is systemic across data-intensive sciences, not unique to ML or medical imaging.

\subsubsection{6.5 Broader Impact}

\textbf{For Policy Makers:} Reproducibility badge programs need quality enforcement, not just presence incentives. Venues could implement:
\begin{itemize}
\item Post-publication audits: Randomly sample badged papers and verify artifact completeness
\item Quality-weighted badges: Distinguish ''Bronze'' (artifact exists) from ''Gold'' (artifact provides detailed specifications)
\item Community-driven quality ratings: Integrate user feedback (GitHub stars, documentation requests) into badge criteria
\end{itemize}

\textbf{For Practitioners:} Our rubric (4 dimensions, 0-10 scale) provides a self-assessment tool. Authors can evaluate their artifacts before submission and iteratively improve quality scores. The low scores for evaluation protocols (1.19/10) and hyperparameters (1.16/10) highlight where to focus effort.

\textbf{For Researchers:} Performance variance (CV) is a scalable reproducibility proxy. While our null result (p=0.418) does not support the specific hypothesis that artifacts reduce variance, the method (aggregating variance from independent attempts) is sound. Future work can apply this approach to other interventions (e.g., preregistration, registered reports, code review).

\subsubsection{6.6 What Would Change the Conclusion?}

Our conclusion---that artifacts exist but quality is insufficient---would be affected under the following conditions:

\begin{enumerate}
\item \textbf{Larger sample (n=100) finds significant effect:} If variance reduction becomes significant at n=100, the current null result is Type II error. The directional trend (d=0.464) makes this plausible.
\end{enumerate}

\begin{enumerate}
\item \textbf{Stratification by domain reveals effects:} Computer vision and NLP may differ in artifact practices. If domain-stratified analysis finds significant effects within CV or NLP, the aggregate null result masks heterogeneity.
\end{enumerate}

\begin{enumerate}
\item \textbf{Artifact quality (not count) predicts variance:} If regression using quality scores (h-m1) instead of binary presence (h-m3) shows negative correlation, it confirms that quality dominates quantity.
\end{enumerate}

None of these possibilities invalidate our core finding: current artifact quality is low. Even if larger samples detect weak effects, the fact remains that mean quality (2.43/10) is far below actionable thresholds.
